% Chapter-specific extensions injected by \almanachinputchapter.
% Keep each extension compact, current, and tied to an engineering decision.

\expandafter\def\csname almanachdeepdive@orientation\endcsname{%
\section{Deep Dive: The Foundation-Model Lifecycle}
\begin{tcbitemize}[skin=sectionraster]
\tcbitem[title={One Model, Many System Layers},raster multicolumn=3]
A foundation model is only one component in a lifecycle that includes data
rights, pretraining, post-training, evaluation, retrieval, tools, serving,
monitoring, and incident response.
The same weights can behave very differently under another prompt template,
retrieval corpus, decoding policy, tool permission, or user workflow.

\tcbitem[title={Capability Is Not a Deployment Decision},raster multicolumn=3]
Translate a proposed capability into a user, decision, error cost, operating
boundary, baseline, acceptance test, fallback, and owner.
Multimodal input or fluent generation expands the interface; it does not remove
the need to define what counts as correct and safe.

\tcbitem[title={Mini Design Review},raster multicolumn=6]
For one AI application, draw
\emph{data $\rightarrow$ model $\rightarrow$ system $\rightarrow$ human
$\rightarrow$ outcome}.
At every arrow, name what can be lost, distorted, attacked, or misunderstood.
Then choose one metric and one qualitative investigation that could reveal the
failure before deployment.
\end{tcbitemize}
}

\expandafter\def\csname almanachdeepdive@mathematics\endcsname{%
\section{Deep Dive: Numerical Stability and Honest Uncertainty}
\begin{tcbitemize}[skin=sectionraster]
\tcbitem[title={Conditioning Before Optimisation},raster multicolumn=3]
An algorithm can be mathematically correct and numerically unreliable.
Inspect scale, units, rank, and condition number before interpreting a large
gradient or unstable coefficient.
Standardisation, orthogonal decompositions, and regularisation often improve
the problem being solved rather than merely speeding up a solver.

\tcbitem[title={Log-Sum-Exp},raster multicolumn=3]
Directly computing $\log\sum_i e^{x_i}$ may overflow.
With $m=\max_i x_i$,
\begin{equation}
  \log\sum_i e^{x_i}=m+\log\sum_i e^{x_i-m}.
\end{equation}
Subtracting $m$ changes neither the result nor softmax ratios, while keeping
exponents in a representable range.

\tcbitem[title={Intervals Need Assumptions},raster multicolumn=6]
A confidence, credible, bootstrap, prediction, or conformal interval answers a
different question under different assumptions.
For every interval, state the random object, conditioning information,
coverage target, data-dependence assumptions, and whether tuning reused the
calibration data.
Check empirical coverage on relevant slices rather than trusting the nominal
percentage by name.
\end{tcbitemize}
}

\expandafter\def\csname almanachdeepdive@programming\endcsname{%
\section{Deep Dive: Reproducible ML Software}
\begin{tcbitemize}[skin=sectionraster]
\tcbitem[title={Notebook to Package},raster multicolumn=3]
Keep exploration in notebooks, but move reusable transformations, models, and
metrics into typed modules with documented interfaces.
A command-line entry point plus a versioned configuration makes the experiment
rerunnable without executing cells in a remembered order.

\tcbitem[title={Test Properties, Not Only Examples},raster multicolumn=3]
Example tests check known cases.
Property and metamorphic tests check invariants: shuffling rows should not
change a permutation-invariant aggregate; serialization should preserve a
prediction within tolerance; a monotone transform should preserve ranks.
These tests catch families of failures that one golden input misses.

\tcbitem[title={Operational Skeleton},raster multicolumn=6]
Build a minimal repository with locked dependencies, configuration validation,
structured logging, data-schema checks, unit tests, one end-to-end smoke test,
an environment capture, and a single command that trains and evaluates a
baseline.
Record code revision, data identifier, seed, hardware, and output artifact in
the same run manifest.
\end{tcbitemize}
}

\expandafter\def\csname almanachdeepdive@data-analytics\endcsname{%
\section{Deep Dive: Data Contracts, Lineage, and Drift}
\begin{tcbitemize}[skin=sectionraster]
\tcbitem[title={A Schema Is Not a Data Contract},raster multicolumn=3]
A contract includes semantics, units, allowed values, null meaning, freshness,
uniqueness, privacy classification, owners, change policy, and service level.
Two columns can share a type while representing incompatible populations or
measurement procedures.

\tcbitem[title={Lineage Makes Errors Reversible},raster multicolumn=3]
Record source, transformation version, validation result, and downstream
consumers for every material dataset and feature.
Lineage should answer which models used a corrupted source and which reports
must be recomputed after a correction.

\tcbitem[title={Drift Is a Question, Not One Score},raster multicolumn=6]
Distinguish input drift $p(x)$, label drift $p(y)$, concept drift $p(y\mid x)$,
and operational changes in collection or policy.
A PSI or KS alert can indicate distribution change but cannot identify cause or
business impact by itself.
Pair statistical signals with slice-level performance, data-quality checks,
and an investigation playbook.
\end{tcbitemize}
}

\expandafter\def\csname almanachdeepdive@machine-learning\endcsname{%
\section{Deep Dive: From Score to Decision Policy}
\begin{tcbitemize}[skin=sectionraster]
\tcbitem[title={Probability, Threshold, Action},raster multicolumn=3]
A classifier score becomes useful only through a decision policy.
Choose thresholds from false-positive and false-negative consequences,
capacity constraints, calibration, and an abstention or review path.
Accuracy at the default threshold rarely captures this operating decision.

\tcbitem[title={Calibration and Coverage},raster multicolumn=3]
Measure reliability curves and proper scoring rules on held-out data.
When uncertainty is material, compare calibrated probabilities, prediction
intervals, or conformal sets and report empirical coverage and set size by
slice.
Coverage without useful sharpness can still produce an unusable system.

\tcbitem[title={Data-Centric Error Loop},raster multicolumn=6]
After a baseline, cluster errors by mechanism: ambiguous target, label error,
missing feature, distribution gap, shortcut, capacity limit, or threshold
policy.
Make one intervention at a time, preserve a protected evaluation set, and
measure whether improvement transfers beyond the error slice that inspired the
change.
\end{tcbitemize}
}

\expandafter\def\csname almanachdeepdive@deep-learning\endcsname{%
\section{Deep Dive: Modern Training and Inference}
\begin{tcbitemize}[skin=sectionraster]
\tcbitem[title={Residual Pre-Norm Blocks},raster multicolumn=3]
Modern deep networks stabilise optimisation with residual paths and
normalisation.
In a pre-norm Transformer block,
$x' = x + \mathrm{Attention}(\mathrm{Norm}(x))$ and
$y = x' + \mathrm{MLP}(\mathrm{Norm}(x'))$.
The residual path carries information and gradients even when a learned branch
is initially poor.

\tcbitem[title={Training Parallelism vs Inference},raster multicolumn=3]
Self-attention can process training positions in parallel, but dense attention
still has quadratic interaction cost in sequence length.
Autoregressive generation remains sequential across new tokens.
A key--value cache avoids recomputing earlier keys and values, trading memory
for lower per-token latency.

\tcbitem[title={Efficiency Is Multi-Dimensional},raster multicolumn=6]
Quantisation, distillation, sparsity, low-rank adaptation, activation
checkpointing, batching, and distributed sharding reduce different costs and
introduce different errors.
Benchmark end-to-end quality, latency, throughput, peak memory, energy, and
operational complexity under the intended workload; parameter count alone is
not a deployment metric.
\end{tcbitemize}
}

\expandafter\def\csname almanachdeepdive@evaluation-mlops\endcsname{%
\section{Deep Dive: Evaluating Generative and Adaptive Systems}
\begin{tcbitemize}[skin=sectionraster]
\tcbitem[title={Quality Is a Vector},raster multicolumn=3]
For generative systems, separate task success, factual grounding, instruction
following, safety, style, latency, and cost.
Aggregate scores can hide a release-blocking failure, so define non-compensable
gates for critical behaviours.

\tcbitem[title={Judges Need Evaluation Too},raster multicolumn=3]
An LLM judge is a measurement instrument.
Calibrate it against blinded human decisions, randomise order, test position
and verbosity bias, report agreement and uncertainty, and keep a stable anchor
set across judge changes.

\tcbitem[title={Sequential Release Evidence},raster multicolumn=6]
Combine offline suites with shadow traffic, canaries, guarded online
experiments, rollback triggers, and post-release monitoring.
Account for repeated looks at online metrics and delayed outcomes.
A release is an evidence update, not permission to retire the evaluation set.
\end{tcbitemize}
}

\expandafter\def\csname almanachdeepdive@reinforcement-learning\endcsname{%
\section{Deep Dive: Offline, Safe, and Preference-Based Learning}
\begin{tcbitemize}[skin=sectionraster]
\tcbitem[title={Offline RL and Coverage},raster multicolumn=3]
Offline RL learns from a fixed behaviour dataset.
Actions outside that dataset's support can receive optimistic values because
their outcomes are weakly identified.
Report behaviour-policy coverage and compare with behaviour cloning and a
conservative non-learning baseline.

\tcbitem[title={Safety Is a Constraint},raster multicolumn=3]
Reward penalties do not guarantee a hard safety limit.
Constrained policies, shields, action filters, simulators, and independent
supervisors address different layers.
Evaluate rare violations, recovery, and reward hacking rather than only mean
return.

\tcbitem[title={Off-Policy Evidence},raster multicolumn=6]
Before deploying a new policy, estimate its value with importance weighting,
doubly robust estimators, or a validated simulator.
Large policy shifts create high-variance or biased estimates when action
probabilities have little overlap.
Report effective sample size, sensitivity to clipping, and where evaluation
depends on extrapolation.
\end{tcbitemize}
}

\expandafter\def\csname almanachdeepdive@trustworthy-ai\endcsname{%
\section{Deep Dive: Agentic Systems and Deterministic Boundaries}
\begin{tcbitemize}[skin=sectionraster]
\tcbitem[title={Treat Tool Output as Untrusted},raster multicolumn=3]
Retrieved documents, web pages, messages, and tool results can contain
instructions designed to compete with system policy.
Parse data with schemas, isolate it from instructions, and keep authorization
outside the model.

\tcbitem[title={Least-Privilege Tool Use},raster multicolumn=3]
Issue short-lived credentials for the minimum resource and action.
Validate arguments, enforce limits, require confirmation for consequential
changes, and log the request, decision, tool result, and artifact versions.
The model proposes; a deterministic boundary disposes.

\tcbitem[title={Trajectory-Level Assurance},raster multicolumn=6]
Evaluate complete trajectories, including retries, tool selection, state
changes, recovery, and human handover.
A component that passes isolated prompts can still fail through compounding
errors, excessive autonomy, or a permission path that bypasses the intended
control.
\end{tcbitemize}
}

\expandafter\def\csname almanachdeepdive@computer-vision\endcsname{%
\section{Deep Dive: Foundation Vision and Robust Perception}
\begin{tcbitemize}[skin=sectionraster]
\tcbitem[title={Prompts Change the Vision Interface},raster multicolumn=3]
Promptable segmentation and vision--language models replace a fixed class list
with points, boxes, text, or examples.
This flexibility expands evaluation: test localisation, open-vocabulary
recognition, grounding, refusal, and consistency across prompt wording.

\tcbitem[title={Shift Is Often Physical},raster multicolumn=3]
Weather, optics, compression, sensor ageing, viewpoint, motion blur, and scene
composition alter the image-generating process.
Corruption benchmarks are useful probes but do not replace data collected in
the intended operational domain.

\tcbitem[title={Perception-to-Action Trace},raster multicolumn=6]
For one detected object, trace calibration, preprocessing, model output,
tracking, fusion, uncertainty, planning, and fallback.
Express how spatial or temporal error propagates into an action margin.
Evaluate the system-level consequence, not only IoU or mAP in isolation.
\end{tcbitemize}
}

\expandafter\def\csname almanachdeepdive@nlp-generative-ai\endcsname{%
\section{Deep Dive: Retrieval, Tools, and Structured Generation}
\begin{tcbitemize}[skin=sectionraster]
\tcbitem[title={RAG Has Two Evaluations},raster multicolumn=3]
Evaluate retrieval with relevance, recall, ranking, freshness, permissions, and
coverage.
Evaluate generation with grounding, citation correctness, completeness, and
task success.
A faithful answer to irrelevant evidence is still a retrieval failure.

\tcbitem[title={Structured Output Is a Contract},raster multicolumn=3]
Define a schema, validate every field, reject unknown actions, and separate
syntactic validity from semantic authorization.
Repair loops need bounded retries and telemetry; silently accepting a plausible
but invalid object moves uncertainty into downstream code.

\tcbitem[title={Long Context Is Not Perfect Memory},raster multicolumn=6]
More tokens increase cost and can dilute relevant evidence.
Compare long-context prompting with retrieval, summarisation, and explicit
state.
Test information at different positions, conflicting evidence, multilingual
inputs, and instructions embedded inside retrieved content.
\end{tcbitemize}
}

\expandafter\def\csname almanachdeepdive@practical-llm-engineering\endcsname{%
\section{Deep Dive: Serving the Adapted Model}
\begin{tcbitemize}[skin=sectionraster]
\tcbitem[title={Quantisation Experiment},raster multicolumn=3]
Compare the reference model with 8-bit and 4-bit weight variants on the same
quality, latency, throughput, and peak-memory suite.
Separate quantised loading from quantisation-aware training and record which
layers or tensors remain at higher precision.

\tcbitem[title={Measure the Workload Shape},raster multicolumn=3]
Prompt length, generated length, batch concurrency, cache reuse, and hardware
determine serving behaviour.
Report time to first token and inter-token latency in addition to total
throughput; an average can hide an unusable tail.

\tcbitem[title={Adapter Governance},raster multicolumn=6]
Version base model, tokenizer, chat template, adapter, merge procedure,
training data, licence, and evaluation result as one deployable bill of
materials.
Verify a merged artifact against the separate adapter path before promotion and
retain a rollback target.
\end{tcbitemize}
}

\expandafter\def\csname almanachdeepdive@fintech\endcsname{%
\section{Deep Dive: Point-in-Time Financial Evaluation}
\begin{tcbitemize}[skin=sectionraster]
\tcbitem[title={What Was Knowable Then?},raster multicolumn=3]
Reconstruct features, universe membership, prices, and filings as they were
available at the decision timestamp.
Later corrections, surviving firms, or revised labels create look-ahead and
survivorship bias even when rows are split by date.

\tcbitem[title={Backtest the Executable Policy},raster multicolumn=3]
Include transaction costs, spread, latency, market impact, turnover,
constraints, rejected orders, and capital limits.
Compare against a simple implementable strategy and test performance across
regimes rather than selecting a convenient interval.

\tcbitem[title={Decisions Need Reasons and Appeals},raster multicolumn=6]
For credit, fraud, or AML workflows, connect model output to threshold,
capacity, adverse-action explanation, human review, and appeal.
Measure calibration and error costs by relevant groups, plus queue burden and
time to resolution.
\end{tcbitemize}
}

\expandafter\def\csname almanachdeepdive@healthcare\endcsname{%
\section{Deep Dive: From Retrospective Accuracy to Clinical Utility}
\begin{tcbitemize}[skin=sectionraster]
\tcbitem[title={External and Prospective Evidence},raster multicolumn=3]
Validate across sites, devices, time periods, workflows, and patient groups.
A retrospective external set tests transportability; a prospective silent
study tests current data flow; an impact study tests whether use changes care
and outcomes.

\tcbitem[title={Thresholds Depend on the Pathway},raster multicolumn=3]
Sensitivity and specificity do not determine value without prevalence,
capacity, downstream tests, delay, and harm.
Decision-curve or utility analysis makes the intervention threshold explicit
and should be complemented by calibration and subgroup uncertainty.

\tcbitem[title={Change Control for Learning Systems},raster multicolumn=6]
Define which data, preprocessing, model, threshold, interface, or workflow
changes require revalidation.
Monitor input and outcome delay, override patterns, incidents, and calibration.
A safe update plan includes rollback and communication, not only a newer model
file.
\end{tcbitemize}
}

\expandafter\def\csname almanachdeepdive@industrial-ai\endcsname{%
\section{Deep Dive: Digital Twins, Edge AI, and Safe Updates}
\begin{tcbitemize}[skin=sectionraster]
\tcbitem[title={Calibrate the Twin},raster multicolumn=3]
A digital twin is useful only within a validated envelope.
Estimate parameters from measured behaviour, compare residuals across operating
regimes, quantify uncertainty, and identify effects absent from the simulator.

\tcbitem[title={Edge Constraints Shape the Model},raster multicolumn=3]
Timing jitter, memory, power, connectivity, sensor bandwidth, and thermal limits
can dominate model choice.
Benchmark on the target device with realistic preprocessing and concurrent
loads, not only on a workstation.

\tcbitem[title={Safe Over-the-Air Change},raster multicolumn=6]
Sign artifacts, bind model and configuration versions, stage rollout, verify a
shadow or canary, protect compatibility, and keep an independently tested
fallback.
An update must preserve the operational safety envelope even when its average
prediction metric improves.
\end{tcbitemize}
}

\expandafter\def\csname almanachdeepdive@supply-chain\endcsname{%
\section{Deep Dive: Forecast Distributions to Robust Decisions}
\begin{tcbitemize}[skin=sectionraster]
\tcbitem[title={Evaluate at the Decision Horizon},raster multicolumn=3]
Use rolling-origin evaluation with the forecast horizons, hierarchy, and data
latency of the planning process.
Report scale-appropriate errors and interval coverage; one aggregate point
metric can hide poor service for intermittent or high-value items.

\tcbitem[title={Reconcile the Hierarchy},raster multicolumn=3]
Product, region, channel, and total forecasts should add consistently.
Reconciliation trades local fit for coherent planning and must respect which
levels contain reliable signal and which are management aggregates.

\tcbitem[title={Optimise Across Scenarios},raster multicolumn=6]
Pass a forecast distribution or scenarios into inventory, routing, or capacity
decisions rather than treating a point forecast as certain demand.
Stress lead time, supplier failure, correlated disruption, and policy delay.
Compare the learned policy with a transparent robust or stochastic-optimisation
baseline.
\end{tcbitemize}
}

\expandafter\def\csname almanachdeepdive@multi-agent-systems\endcsname{%
\section{Deep Dive: LLM Agents as Distributed Systems}
\begin{tcbitemize}[skin=sectionraster]
\tcbitem[title={Distinguish Three Agent Traditions},raster multicolumn=3]
Classical BDI agents execute explicit beliefs, desires, and intentions.
MARL agents learn policies through interaction.
LLM agents generate plans and tool calls from context.
The three can be combined, but their state, guarantees, and failure modes are
not interchangeable.

\tcbitem[title={Protocols Need Semantics},raster multicolumn=3]
A message schema states syntax; an interaction protocol also defines allowed
states, idempotency, timeouts, authorization, conflict resolution, and recovery.
Use correlation identifiers and explicit ownership so retries do not create
duplicate consequential actions.

\tcbitem[title={Evaluate Trajectories and Emergence},raster multicolumn=6]
Measure task success, steps, tool errors, communication cost, latency,
permission violations, deadlocks, loops, and human interventions.
Test adversarial messages and collusion as well as cooperative examples.
Local pass rates do not prove stable collective behaviour.
\end{tcbitemize}
}

\expandafter\def\csname almanachdeepdive@commerce-marketing\endcsname{%
\section{Deep Dive: Feedback Loops and Multi-Sided Objectives}
\begin{tcbitemize}[skin=sectionraster]
\tcbitem[title={Logged Data Reflects the Old Policy},raster multicolumn=3]
Clicks and purchases are observed only for items exposed by an earlier ranking
and interface.
Position, availability, and selection bias confound naive supervised targets.
Randomised exploration or carefully justified propensity methods help estimate
counterfactual performance.

\tcbitem[title={Delayed and Interfering Outcomes},raster multicolumn=3]
Campaigns can change later conversion, returns, fatigue, and word of mouth.
One user's exposure can affect inventory, prices, sellers, and other users, so
standard independent-unit A/B assumptions may fail.

\tcbitem[title={Balance the Marketplace},raster multicolumn=6]
Evaluate relevance, conversion, margin, diversity, novelty, long-term retention,
producer exposure, support burden, and safety constraints.
Define which objectives are hard constraints and which can trade off.
Monitor feedback loops that make the training data increasingly reflect the
model's own past choices.
\end{tcbitemize}
}

\expandafter\def\csname almanachdeepdive@it-law\endcsname{%
\section{Deep Dive: From Legal Classification to Engineering Evidence}
\begin{tcbitemize}[skin=sectionraster]
\tcbitem[title={Begin with Role and Intended Purpose},raster multicolumn=3]
Map jurisdictions, actor roles, intended purpose, affected people, sector,
data flows, and lifecycle stage before selecting obligations.
Provider, deployer, importer, distributor, employer, controller, and processor
can carry different duties for the same technical component.

\tcbitem[title={Translate Duties into Controls},raster multicolumn=3]
Turn requirements into owners, controls, evidence, retention, review dates, and
incident paths.
Examples include provenance records, risk management, human oversight,
cybersecurity, performance monitoring, notices, rights handling, and supplier
contracts.

\tcbitem[title={Maintain a Legal Change Log},raster multicolumn=6]
Record the source, jurisdiction, version, effective date, interpretation owner,
affected system, and next review.
The legal chapter is orientation, not legal advice; current official text and
qualified counsel govern consequential decisions.
Reassess when purpose, data, model, geography, users, or autonomy changes.
\end{tcbitemize}
}

\expandafter\def\csname almanachdeepdive@visualization\endcsname{%
\section{Deep Dive: Uncertainty, Accessibility, and Monitoring}
\begin{tcbitemize}[skin=sectionraster]
\tcbitem[title={Show the Distribution Behind the Mean},raster multicolumn=3]
Use intervals, densities, small multiples, or raw observations to show spread,
sample size, and heterogeneity.
For probabilistic models, reliability diagrams and risk--coverage curves reveal
behaviour hidden by a single discrimination score.

\tcbitem[title={Design Beyond Colour},raster multicolumn=3]
Use contrast, direct labels, position, shape, line style, and ordering so the
message survives colour-vision differences and monochrome output.
Provide a meaningful caption or alternate description and preserve logical
reading order in the exported document.

\tcbitem[title={A Monitoring Chart Must Trigger Action},raster multicolumn=6]
Connect every panel to an owner, threshold, investigation, and time horizon.
Separate service health, data quality, prediction distribution, delayed
performance, subgroup behaviour, incidents, and business outcomes.
Annotate model or policy changes so the chart can support causal investigation
rather than merely display motion.
\end{tcbitemize}
}

\expandafter\def\csname almanachdeepdive@startup\endcsname{%
\section{Deep Dive: Portfolio Evidence and Option Value}
\begin{tcbitemize}[skin=sectionraster]
\tcbitem[title={Fund Learning Milestones},raster multicolumn=3]
Tie spending to the next uncertainty-reducing milestone: verified problem,
representative data, baseline feasibility, paid workflow adoption, responsible
operation, or repeatable economics.
Do not confuse accumulated code with reduced venture risk.

\tcbitem[title={Preserve Strategic Options},raster multicolumn=3]
Design pilots so data, evaluations, interfaces, and customer learning remain
useful if the model, supplier, or route to market changes.
Portability and a credible manual fallback create option value when technology
and regulation move faster than the plan.
\end{tcbitemize}
}

\expandafter\def\csname almanachdeepdive@scientific-writing\endcsname{%
\section{Deep Dive: Reproducibility Audit}
\begin{tcbitemize}[skin=sectionraster]
\tcbitem[title={Recompute from a Clean Environment},raster multicolumn=3]
Give the repository and manuscript to a reader who did not create the original
environment.
Record missing data, secrets, undocumented preprocessing, manual steps,
unstable downloads, and discrepancies between figures and reported numbers.

\tcbitem[title={Audit the Claim Graph},raster multicolumn=3]
For every abstract and conclusion claim, link the supporting result, method,
artifact, and limitation.
Remove orphan claims and propagate changed numbers or caveats through captions,
tables, discussion, model cards, and supplementary material.
\end{tcbitemize}
}

\newcommand{\almanachdeepdive}[1]{%
    \ifcsname almanachdeepdive@#1\endcsname
        \csname almanachdeepdive@#1\endcsname
    \fi
}
